% Shared preamble for the technical reference manual.
% Authored file - safe to edit. Regenerating the documentation never touches it.

\usepackage[T1]{fontenc}
\usepackage[utf8]{inputenc}
\usepackage{lmodern}
\usepackage{microtype}

\usepackage[
  a4paper,
  top=24mm,
  bottom=24mm,
  left=23mm,
  right=23mm,
  headheight=14pt
]{geometry}

% The contents lists sections (one per documented object) but not the method
% subsections nested inside class entries; those are reachable through the
% class entry, the object indexes, and the PDF bookmarks.
\setcounter{tocdepth}{1}

\usepackage{amsmath}
\usepackage{amssymb}
\usepackage{array}
\usepackage{booktabs}
\usepackage{longtable}
\usepackage{enumitem}
\usepackage{xcolor}
\usepackage{listings}
\usepackage{graphicx}
\usepackage{fancyhdr}
\usepackage[bottom]{footmisc}
\usepackage{makeidx}

\usepackage{multicol}

\usepackage{tikz}
\usetikzlibrary{arrows.meta, positioning, fit, calc, shapes.geometric, backgrounds}

\usepackage[
  colorlinks=true,
  linkcolor=black,
  citecolor=black,
  urlcolor=black,
  bookmarksnumbered=true,
  bookmarksopen=true,
  pdfborder={0 0 0}
]{hyperref}
\usepackage[nameinlink]{cleveref}

% ---------------------------------------------------------------------------
% Monochrome-friendly palette. Every figure and listing must remain legible
% when printed in greyscale, so only neutral tones are defined.
% ---------------------------------------------------------------------------
\definecolor{rulegrey}{gray}{0.45}
\definecolor{shadegrey}{gray}{0.94}
\definecolor{codegrey}{gray}{0.30}

% ---------------------------------------------------------------------------
% Compact sectioning.
%
% The reference part contains one section per documented object. The class
% defaults put roughly a centimetre of vertical space around each heading,
% which at that density costs well over a hundred pages on its own. These are
% the standard \@startsection parameters with the skips tightened; no
% additional package is required.
% ---------------------------------------------------------------------------
\makeatletter
\renewcommand\section{\@startsection{section}{1}{\z@}%
  {-2.2ex \@plus -0.6ex \@minus -.2ex}%
  {0.7ex \@plus .2ex}%
  {\normalfont\large\bfseries\raggedright}}
\renewcommand\subsection{\@startsection{subsection}{2}{\z@}%
  {-1.6ex \@plus -0.5ex \@minus -.2ex}%
  {0.4ex \@plus .2ex}%
  {\normalfont\normalsize\bfseries\raggedright}}
\renewcommand\subsubsection{\@startsection{subsubsection}{3}{\z@}%
  {-1.2ex \@plus -0.4ex \@minus -.2ex}%
  {0.2ex \@plus .1ex}%
  {\normalfont\normalsize\itshape\raggedright}}
\makeatother

% ---------------------------------------------------------------------------
% Reference-entry layout, modelled on the reference manuals CRAN generates for
% R packages: a run-in bold field label followed immediately by its content,
% rather than a numbered subsection per field.
% ---------------------------------------------------------------------------

% Run-in field label: "Usage:", "Arguments:", "Value:", ...
\newcommand{\apifield}[1]{%
  \par\addvspace{2.5pt}\noindent\textbf{#1:}\enspace\ignorespaces}

% Hanging-indent list used for arguments, fields, and enumeration members.
\newenvironment{arglist}%
  {\begin{description}[
      style=sameline,
      leftmargin=2.2em,
      labelindent=0pt,
      labelsep=0.5em,
      topsep=2pt,
      itemsep=1pt,
      parsep=0pt,
      font=\normalfont]}%
  {\end{description}}

% Keeps a short entry from being split across a page break at its heading.
\newcommand{\apibreak}{\filbreak}

% ---------------------------------------------------------------------------
% Source listings.
% ---------------------------------------------------------------------------
\lstset{
  basicstyle=\ttfamily\small,
  keywordstyle=\ttfamily\bfseries,
  commentstyle=\ttfamily\itshape\color{codegrey},
  stringstyle=\ttfamily,
  showstringspaces=false,
  breaklines=true,
  breakatwhitespace=false,
  columns=fullflexible,
  keepspaces=true,
  frame=single,
  rulecolor=\color{rulegrey},
  framesep=4pt,
  backgroundcolor=\color{shadegrey},
  xleftmargin=6pt,
  xrightmargin=2pt,
  aboveskip=8pt,
  belowskip=8pt,
  literate={-}{{-}}1 {_}{{\_}}1
}

% Signature blocks in the reference entries: unframed, indented, and tight,
% so a one-line signature costs one line of output.
\lstdefinestyle{usage}{
  basicstyle=\ttfamily\small,
  frame=none,
  backgroundcolor=\color{white},
  xleftmargin=1.6em,
  xrightmargin=0pt,
  aboveskip=1.5pt,
  belowskip=2.5pt,
  breaklines=true,
  breakindent=2em,
  postbreak={},
  columns=fullflexible
}

% ---------------------------------------------------------------------------
% Table defaults. Generated tables use longtable + booktabs throughout.
% ---------------------------------------------------------------------------
\setlength{\tabcolsep}{4pt}
\renewcommand{\arraystretch}{1.05}
\setlength{\LTcapwidth}{\textwidth}
\raggedbottom

% Generated tables carry long typewriter identifiers. TeX does not hyphenate
% typewriter text, so give the paragraph builder room to stretch rather than
% overflowing the column; the generator additionally inserts \allowbreak
% opportunities inside qualified names.
\setlength{\emergencystretch}{3em}
\tolerance=2000
\hbadness=4000
\hfuzz=2pt

% ---------------------------------------------------------------------------
% Page style.
% ---------------------------------------------------------------------------
\pagestyle{fancy}
\fancyhf{}
\fancyhead[L]{\small\nouppercase{\leftmark}}
\fancyhead[R]{\small\thepage}
\fancyfoot[L]{\small\projectname\ \projectversion}
\fancyfoot[R]{\small rev.\ \texttt{\projectrevisionshort}}
\renewcommand{\headrulewidth}{0.4pt}
\renewcommand{\footrulewidth}{0.4pt}

\fancypagestyle{plain}{%
  \fancyhf{}%
  \fancyfoot[L]{\small\projectname\ \projectversion}%
  \fancyfoot[R]{\small rev.\ \texttt{\projectrevisionshort}}%
  \renewcommand{\headrulewidth}{0pt}%
  \renewcommand{\footrulewidth}{0.4pt}%
}

% ---------------------------------------------------------------------------
% Cross-reference naming for the object kinds this manual defines.
% ---------------------------------------------------------------------------
\crefname{section}{section}{sections}
\Crefname{section}{Section}{Sections}
\crefname{chapter}{chapter}{chapters}
\Crefname{chapter}{Chapter}{Chapters}
\crefname{figure}{figure}{figures}
\Crefname{figure}{Figure}{Figures}
\crefname{table}{table}{tables}
\Crefname{table}{Table}{Tables}

% ---------------------------------------------------------------------------
% Shared semantic macros.
% ---------------------------------------------------------------------------
\newcommand{\pkg}[1]{\texttt{#1}}
\newcommand{\srcfile}[1]{\texttt{#1}}
\newcommand{\envvar}[1]{\texttt{#1}}

% Marks a statement that could not be established from repository evidence.
\newcommand{\noevidence}{%
  \par\medskip\noindent\fbox{\parbox{0.96\linewidth}{%
    \small\textbf{No repository evidence.} Rationale cannot be determined from
    repository evidence.}}\par\medskip%
}

\makeindex
